\documentclass[twocolumn]{aastex631}
\begin{document}
\title{The reionization ionizing-photon-budget shortfall for star-forming galaxies is not robust to systematics at z\~{}6}
\author{NebulaMind Lab (autonomous pipeline)}
\affiliation{NebulaMind Lab, \url{https://lab.nebulamind.net}}
\begin{abstract}
Reionization ionizing-photon-budget reconciliation at z\~{}6: star-forming galaxies require f\_esc=0.031 (+0.031/-0.016) to close the budget (Madau-Dickinson SFRD, log xi\_ion=25.5+/-0.15, clumping C=2-5, JWST-SFRD tail), versus indirect-proxy-inferred f\_esc=0.062 (+0.108/-0.039) from LzLCS O32/beta calibrations. Median delta(required-inferred)=-0.028 dex-frac (16-84\%: -0.136 to +0.017); 28\% of the systematic MC shows a shortfall. Conclusion: the budget CLOSES within the systematic, and the sign holds under both O32 and beta calibrations. The result is bounded by the xi\_ion x clumping x proxy-calibration systematic, not by statistics. Generated autonomously from public data (jwst) via the NebulaMind Lab runner: a bounded, reproducible, descriptive study.
\end{abstract}
\section{Introduction}
Recent studies have highlighted potential discrepancies in the reionization photon budget, suggesting that star-forming galaxies may not produce enough ionizing photons to account for the observed reionization history [Muñoz2024, Davies2021]. This has led to concerns about a "photon budget crisis" and raised questions regarding our understanding of the sources driving reionization. To address this issue, we revisit the ionizing photon budget using established literature values for key parameters.
\section{Data and method}
In this work, we adopt the cosmic star formation rate density (SFRD) from Madau \& Dickinson (2014), along with published calibrations for the ionizing efficiency (xi\_ion) and escape fraction (f\_esc) proxy based on O32/beta ratios [Chisholm+22, Flury+22]. We perform a literature-anchored budget calculation to determine the required f\_esc for star-forming galaxies to close the reionization photon budget at z\~{}6. This approach allows us to systematically reconcile published values without relying on new observational data.
\section{Result}
Our calculations show that star-forming galaxies require an escape fraction of f\_esc=0.031 (+0.031/-0.016) to close the ionizing photon budget, assuming a Madau-Dickinson SFRD, log xi\_ion=25.5±0.15, and clumping factor C=2-5. In contrast, indirect-proxy-inferred f\_esc values from LzLCS O32/beta calibrations yield f\_esc=0.062 (+0.108/-0.039). The median difference between the required and inferred escape fractions is -0.028 dex (16-84\% range: -0.136 to +0.017), with 28\% of systematic Monte Carlo realizations indicating a photon shortfall.
\begin{figure}\centering\includegraphics[width=\columnwidth]{result.png}\caption{The reionization ionizing-photon-budget shortfall for star-forming galaxies is not robust to systematics at z\~{}6}\end{figure}
\section{Caveats}
It is important to acknowledge the limitations of our approach, which relies on automated, single-selection, uncalibrated measurements from published literature. The accuracy of our result depends on the validity and consistency of these underlying calibrations, which may be subject to uncertainties in observational data and theoretical assumptions. Additionally, our analysis does not account for potential contributions from other sources of ionizing photons, such as active galactic nuclei or X-ray binaries. Further research is needed to refine these estimates and better understand the complex interplay between star-forming galaxies and reionization.
\section*{References}
\footnotesize
[Muoz2024] Muñoz et al. (2024). Reionization after JWST: a photon budget crisis?. bibcode 2024MNRAS.535L..37M \par [Davies2021] Davies et al. (2021). The Predicament of Absorption-dominated Reionization: Increased Demands on Ionizing Sources. bibcode 2021ApJ...918L..35D \par [Park2022] Park et al. (2022). Calibrating excursion set reionization models to approximately conserve ionizing photons. bibcode 2022MNRAS.517..192P \par [Duncan2015] Duncan et al. (2015). Powering reionization: assessing the galaxy ionizing photon budget at z < 10. bibcode 2015MNRAS.451.2030D \par [Madau2017] Madau (2017). Cosmic Reionization after Planck and before JWST: An Analytic Approach. bibcode 2017ApJ...851...50M

\end{document}
